\section{Magnetar Kind ($\textnormal{kind} = 0x03$)}
\label{sec:magnetar-kind}

\subsection{Slot and ABI}

\texttt{0x012205} with leading \texttt. Same physical
precompile address as Pulsar (kind $0x01$) and Corona
(kind $0x02$); dispatch resolves to the Magnetar
\texttt{P3QVerifier} implementation.

\begin{lstlisting}[language={}]
// EVM precompile at slot 0x012205, kind = 0x03
bool ok = P3Q.verify(
    0x03,                            // Magnetar kind
    magnetarSig,                     // FIPS-205 SLH-DSA threshold sig
    messageHash,                     // H(prev_root || new_root || batch_hash)
    groupPubKey                      // Magnetar group public key
);
\end{lstlisting}

\subsection{Underlying scheme}

Magnetar is the Lux deployment of FIPS-205 SLH-DSA~\cite{fips205}
under a Lux-defined threshold-aggregation wrapper specified in
LP-181~\cite{lp181}. By construction a completed threshold
signature is byte-identical to a single-party FIPS-205 SLH-DSA
signature on the same message and group public key; the verifier
core is therefore the FIPS-205 SLH-DSA verifier with no threshold-
specific logic. The no-reconstruct threshold signer that produces
that signature is research-stage --- full no-reconstruct threshold
SLH-DSA is an open problem --- so the on-chain path stays
fail-closed (below) until it is proven.

SLH-DSA is hash-based: hardness reduces to the collision
resistance of the underlying hash function (SHA-256 or SHA-3
configurations). This is the only PQ family in P3Q whose
hardness does not depend on lattice assumptions; Magnetar is
the cross-family last-line PQ kind.

\subsection{Wire format}

\begin{table}[h]
\centering
\small
\begin{tabular}{@{}rrl@{}}
\toprule
\textbf{Offset} & \textbf{Size} & \textbf{Field} \\
\midrule
0 & 1 & \texttt{kind = 0x03} \\
1 & 1 & \texttt{mode = 0x01} (reserved Magnetar parameter set) \\
2 & 4 & \texttt{bitmapLen} (BE32) \\
6 & \texttt{bitmapLen} & participant bitmap (set bit = participating) \\
$6{+}\textit{blen}$ & 4 & \texttt{sigLen} (BE32) \\
$10{+}\textit{blen}$ & \texttt{sigLen} & \texttt{slhSig} (\textasciitilde\SI{35664}{B} for 192f) \\
$10{+}\textit{blen}{+}\textit{slen}$ & 4 & \texttt{pkLen} (BE32) \\
$14{+}\textit{blen}{+}\textit{slen}$ & \texttt{pkLen} & \texttt{groupPubKey} \\
$14{+}\textit{blen}{+}\textit{slen}{+}\textit{pkLen}$ & 32 & \texttt{messageHash} \\
\bottomrule
\end{tabular}
\caption{Magnetar-kind wire format. The participant bitmap is
specific to the Magnetar kind -- the threshold-aggregation
wrapper needs to know which signers contributed in order to
re-derive the aggregated public key for verification. Pulsar
($0x01$) and Corona ($0x02$) carry that information inside the
sig structure; Magnetar's hash-tree structure puts it in the
wire envelope.}
\label{tab:magnetar-wire}
\end{table}

\subsection{Verifier core}

The Go reference implementation is planned at
\texttt{lux/standard/precompiles/p3q\_family/magnetar.go} and
will wrap \texttt{luxfi/magnetar/verify}. The core is FIPS-205
SLH-DSA conformant (any conformant FIPS-205 verifier suffices).

Status as of 2026-06-04: same as Corona, the verifier core is
not yet wired into the P3Q dispatch table; \texttt{p3q.Run}
routes \texttt to \texttt{abiFalse} via the
\texttt{ErrKindNotImplemented} path. Solidity callers MAY emit
\texttt{P3Q.verify(0x03, ...)} from network genesis; the call
charges the per-kind gas and returns \texttt{false} until the
verifier lands.

\subsection{Parameter set choice: 192f for production
(locked 2026-06-04)}
\label{sec:magnetar-192f-lock}

SLH-DSA at the 192-bit security level has two parameter sets;
the production deployment is locked to the ``fast'' variant
after the 2026-06-04 measurement.

\begin{table}[h]
\centering
\small
\begin{tabular}{@{}lrrrl@{}}
\toprule
\textbf{Parameter} & \textbf{Sig size} & \textbf{Sign (Spark)} & \textbf{Sign (M1 Max)} & \textbf{Production fit} \\
\midrule
SLH-DSA-192s (small) & $\sim$\SI{16}{KB} & \textbf{\SI{795}{\milli\second}} & \SI{1575}{\milli\second} & breaks \SI{1}{\second} round budget \\
\textbf{SLH-DSA-192f (fast)} & $\sim$\SI{35}{KB} & \textbf{\SI{36}{\milli\second}} & \SI{71}{\milli\second} & \textbf{fits LP-205 round budget} \\
\bottomrule
\end{tabular}
\caption{SLH-DSA 192s vs 192f per-validator sign times,
measured in the Quasar 4-scheme matrix bench
(\texttt{/tmp/quasar-matrix-bench}, 2026-06-04, mean of 50
rounds at $N{=}64$). 192f is $22\times$ faster than 192s on
Spark. Source~\cite{parityaudit}.}
\label{tab:magnetar-192f-lock}
\end{table}

The matrix bench measured Magnetar-192s as the wall-clock
bottleneck of every PQ-heavy cell. At $N{=}64$ on Spark,
PQ-heavy total wall-clock with 192s was \SI{1.33}{\second} --
blowing past the \SI{1}{\second} sub-second-finality target. The
bottleneck is the per-validator SLH-DSA sign, not the threshold
aggregation, not the verifier, not the network -- it is the
inherent FORS + WOTS+ hash-tree construction cost on the larger
192s tree.

Switching to 192f drops the per-validator sign from
\SI{795}{\milli\second} to \SI{36}{\milli\second}
($22\times$ faster) at the cost of a $\sim$$2.2\times$ larger
signature ($\sim$\SI{35}{KB} vs $\sim$\SI{16}{KB}). For a
PQ-heavy rollup posting one Magnetar verify per batch, the
extra \SI{19}{KB} on the wire is negligible against the
\SI{0.76}{\second} latency saved.

\subsubsection{Production lock}

LP-222-conformant deployments \textbf{MUST} configure Magnetar
with SLH-DSA-192f. The threshold-init API at
\texttt{lux/threshold/protocols/magnetar/} accepts a parameter-
set argument; PQ-heavy chains pass
\texttt{magnetar.ParamSetFast192}. Operators choosing 192s for
the $2.2\times$ smaller signature opt into multi-second finality
and MUST run \texttt{mode=PQ-strict} (no Magnetar leg) for the
sub-second tier; 192s is reserved for long-archival anchoring
use cases that don't require fast finality.

Verification cost is symmetric between the two sets
($\sim$\SI{1.92}{\milli\second} CPU); the parameter set affects
only the wall-clock of the per-validator signing ceremony, not
the on-chain verify cost. The Magnetar kind's gas charge is
\SI{30000}{gas} regardless of parameter set (signature size on
the wire varies but the verifier work is bounded by the same
hash-tree-walk constant).

\subsection{Gas cost}

\textbf{\SI{30000}{gas}} -- comparable to the Pulsar kind's
\SI{50000}{}. SLH-DSA verify is fast (\textasciitilde\SI{1.92}{\milli\second}
CPU per the parity audit -- comparable to Corona's
\textasciitilde\SI{1.6}{\milli\second} CPU but with very
different gas cost) because the verifier walks a hash tree
rather than performing lattice NTT inverse.

The gas calibration reflects the modest CPU verify cost
relative to Corona; Magnetar verify is dominated by memory
bandwidth on the WOTS+ tree-walk, which is a tight loop on
SHA-256 / SHA-3 hash function evaluation that modern CPUs
execute efficiently.

\subsection{Performance}

\begin{table}[h]
\centering
\small
\begin{tabular}{@{}lrrl@{}}
\toprule
\textbf{Metric} & \textbf{M1 Max} & \textbf{Spark} & \textbf{Source} \\
\midrule
Per-validator sign 192f at $N{=}64$ & \SI{71}{\milli\second} & \SI{36}{\milli\second} & matrix bench (measured)~\cite{parityaudit} \\
Per-validator sign 192s at $N{=}64$ & \SI{1575}{\milli\second} & \SI{795}{\milli\second} & matrix bench (measured) \\
CPU verify (single sig) & $\sim$\SI{1.92}{\milli\second} & $\sim$\SI{1.92}{\milli\second} & parity audit (measured) \\
Blackwell sm\_120 verify & --- & $\sim$\SI{50}{\milli\second} & projected (LP-203 kernel)~\cite{lp203} \\
Blackwell sm\_120 batched ($N{=}64$) & --- & $\sim$\SI{80}{\milli\second} aggregate & projected \\
\bottomrule
\end{tabular}
\caption{Magnetar kind performance. The 192s/192f signing
delta is the dominant operational consideration; verify cost is
symmetric between sets. Blackwell verify timings are
projections.}
\label{tab:magnetar-perf}
\end{table}

\subsection{Use case}

A rollup that wants \textbf{hash-based family independence}
from lattice families configures its sequencer to sign each
batch with a Magnetar threshold sig in addition to Pulsar (and
possibly Corona). The rollup contract calls
\texttt{P3Q.verify(0x01, ...)} then \texttt{P3Q.verify(0x03, ...)}
at slot \texttt{0x012205}.

SLH-DSA security relies only on the collision resistance of its
underlying hash function; it is immune to advances in lattice
cryptanalysis and would survive even a complete break of
Module-LWE (which both Pulsar and Corona rest on). Magnetar is the \textbf{last-line PQ}
-- the kind of choice for sovereign-archival commitments,
long-tenor custody attestations, cross-chain bridges, and
regulator-facing audit trails. The cost (\SI{30000}{gas} per
verify $+$ $\sim$\SI{35}{KB} signature on the wire) is
reasonable for low-frequency, high-value batches.
